% Class-compatibility preamble retained for a reader excerpt of
%   Sahar Diskin, Anna Geisler, "Perfect Matching in Product Graphs and in their
%   Random Subgraphs", arXiv:2404.14020v4, CC BY 4.0.
%
% The source pins its own publisher class (scrartcl); the delivery wrapper accepts
% only the classes article and amsart, so this file supplies the amsart baseline and
% then copies the author's own declaration block mechanically and unchanged
% (pinned-original-source.tex, source lines 30-57: the author's newtheorem
% declarations and the author's macro declarations).  The author's own package
% list of source lines 4-21 is reused only for the entries the retained text needs
% (enumitem, amsmath, amssymb, amsthm, mathtools, xcolor, geometry) and the
% cross-reference packages already needed by the author's own \Cref/\eqref usage
% (hyperref, cleveref); tikz, caption, subcaption, float, graphicx, wrapfig and
% abstract serve figures and the title block that this excerpt does not carry.
% No mathematical macro, statement, symbol or assumption is introduced, changed
% or removed by this file.
\documentclass[10pt,a4paper]{amsart}
\usepackage[shortlabels]{enumitem}
\usepackage{amsmath}
\usepackage{amssymb}
\usepackage{amsthm}
\usepackage{mathtools}
\usepackage{xcolor}
\usepackage[margin=1in]{geometry}
\usepackage{hyperref}
\usepackage{cleveref}
\hypersetup{
  colorlinks=true,
  linkcolor=blue,
  filecolor=magenta,
  urlcolor=cyan
}

\newtheorem{theorem}{Theorem}
\newtheorem{thm}{Theorem}[section]
\newtheorem{corollary}[thm]{Corollary}
\newtheorem{lemma}[thm]{Lemma}
\newtheorem{proposition}[thm]{Proposition}
\newtheorem{claim}[thm]{Claim}
\newtheorem{remark}[thm]{Remark}
\newtheorem{question}[thm]{Question}

\newcommand{\whp}{whp}
\newcommand{\defi}{\operatorname{def}}
\newcommand{\odd}{\operatorname{odd}}
\newcommand{\iso}{\operatorname{iso}}
\newcommand{\cN}{\mathcal{N}}
\newcommand{\cO}{\mathcal{O}}
\newcommand{\cE}{\mathcal{E}}
\newcommand{\cK}{\mathcal{K}}
\newcommand{\cP}{\mathcal{P}}
\newcommand{\cQ}{\mathcal{Q}}

\newcommand{\footremember}[2]{%
  \footnote{#2}%
  \newcounter{#1}%
  \setcounter{#1}{\value{footnote}}%
}

\newcommand{\ag}[1]{\textcolor{purple}{(Anna: #1)}}
\newcommand{\sd}[1]{\textcolor{red}{(Sahar: #1)}}
\title{An exponential upper bound on the number of bipartitions of a sparse-boundary set with few internal edges}
\author{Sahar Diskin, Anna Geisler}
\date{arxiv:2404.14020v4; CC BY 4.0}
\begin{document}
\maketitle
\noindent\textit{Editorial note.} All delivered mathematical text is the author's own, copied from the pinned author source sources/combinatorics-fl04-vspace-1-5cm-perfect-matching-in-p/pinned-original-source.tex; no statement or proof step was rewritten, reordered or removed.\par\medskip
\noindent\textit{Editorial note.} Preamble provenance: the source pins the publisher class scrartcl, which the delivery wrapper does not accept (supported wrapper classes are article and amsart), so this card uses sources/combinatorics-fl04-vspace-1-5cm-perfect-matching-in-p/retained-preamble.tex, which supplies the amsart baseline and then carries the author's own newtheorem and macro declarations (pinned source lines 30-57) verbatim. No macro was invented and no mathematical symbol was changed.\par\medskip
\noindent\textit{Editorial note.} Bibliography provenance: the source itself carries no thebibliography environment and ends with bibliographystyle abbrv plus bibliography perc (author source lines 1691-1692), so the bibliography is resolved with BibTeX; the database used is the identical copy staged with the same arXiv version under sources/product-graph-matchings-2404.14020v4/perc.bib, whose licence evidence is also status supported for source\_id arxiv:2404.14020v4. No entry was edited and the bibliography style is the author's own abbrv.\par\medskip
\noindent\textit{Editorial note.} The author's displayed consequence labelled eq:colouring-weight at source lines 1360-1362 is delivered verbatim inside the statement span, so its label resolves inside the assembled material and no stub is needed for it.\par\medskip
\noindent\textit{Editorial note.} Two displays inside the counted proof carry the author's labels eq:counting-lem and eq:size-X; they are inside the delivered span and are used by the proof itself.\par\medskip
\noindent\textit{Editorial note.} The proof cites nothing, so the card needs no navigation stub; the bibtex database is declared because the source's own bibliography command is the only bibliography the paper has.\par\medskip
\noindent\textit{Allowed standard prerequisites (editorial).}\begin{itemize}
\item The graph G with its dimension-regularity parameters, the constants M and the relation ln n >= (ln 2/(M-1)) d. Set up in the author's introduction and preliminaries (source lines 201-240); used at the end of the counted proof. Supplies the parameter relation that turns the exponential factor 2\textasciicircum{}\{|R|\} (1-p\_-)\textasciicircum{}\{2(M-1)|R|\} into exp(-Omega(|R|)).
\end{itemize}\medskip
\begin{lemma}
\label{lem:two-colour-count}
Fix constants $K, K'>0$.
Let $R\subseteq V(G)$ have size $|R|\ge d$ and
\[
 e_G(R,R^c)\le K |R|\ln d.
\]
The number of bipartitions $R=R_1\sqcup R_2$ satisfying
\[
 e_G(R_1)+e_G(R_2)\le K' |R|
\]
is at most
\[
 \exp\left\{O\left(\frac{|R|\ln^2 d}{d}\right)\right\}.
\]
Consequently, for every fixed $\theta>0$,
\begin{equation}\label{eq:colouring-weight}
 \sum_{R=R_1\sqcup R_2} (1-p_-)^{\theta(e_G(R_1)+e_G(R_2))}
 \le\exp\{-o(|R|)\}.
\end{equation}
\end{lemma}

\begin{proof}
Put $H=G[R]$ and $m=e_G(R_1)+e_G(R_2)$, and assume w.l.o.g. that $|R_2|\ge|R_1|$.
Note that
\begin{align} \label{eq:counting-lem}
   e_G(R, R^c)=d|R|- \sum_{v \in R} d_H(v), \qquad 2m= \sum_{v \in R_1} d_{H[R_1]} (v) + \sum_{v \in R_2} d_{H[R_2]}(v). 
\end{align}
Let $X$ be the set of vertices $v \in R$ with either few neighbours in $H$ or many neighbours in their own bipartition class, i.e.,
\[
 X\coloneqq \left\{v \in R: d_H(v)<\frac{3}{4}d \qquad\text{or}\qquad  d_{H[R_i]}(v)>d/4\quad( \text{ if } v\in R_i)\right\}.
\]
By definition every $v \in R\setminus V$ has at least $\frac{d}{2}$ neighbours in the opposite bipartition class.
From \eqref{eq:counting-lem} we conclude that
\begin{align} \label{eq:size-X}
    |X| \leq \frac{4 e_G(R, R^c)}{d} + \frac{8m}{d}.
\end{align}

Form a random set $T$ by choosing each vertex of $R_2$ independently with probability $\rho=12\ln d/d$. The expected number of vertices of $R_1\setminus X$ with no neighbour in $T$ is at most
\[
 s(1-\rho)^{d/2}\le \frac{|R|}{d^6}.
\]
Moreover,
\[
 \mathbb{E}|T|= \rho |R_2| \le\rho |R|,
 \qquad
 \mathbb{E}|R_2\cap N_H(T)|\le\rho 2 e_G(R_2)\le2\rho m.
\]
Applying Markov's inequality to these three quantities gives a set $T\subseteq R_2$ with
\[
 |T|\le3\rho |R|,
 \qquad
 |R_1\setminus N_H(T)|\le |X| + \frac{|R|}{d^2},
 %\leq \frac{4b(R)}d+\frac{8m}d+\frac{s}{d^2},
 \qquad
 |R_2\cap N_H(T)|\le6\rho m.
\]
Combining \eqref{eq:size-X} with our assumptions on $e_G(R, R^c)$ and $m$, we obtain that $R_1\mathbin{\triangle}N_H(T)$ has size $O(|R|\ln d/d)$.
The pair $(T, R_1\mathbin{\triangle}N_H(T))$ determines $R_1$, and the number of possible pairs is at most
\[
 \left(\sum_{j \le O(|R|\ln d/d)}\binom{|R|}{j}\right)^2
 =\exp\left\{O\left(\frac{|R|\ln^2 d}{d}\right)\right\},
\]
which proves the first assertion.

For \eqref{eq:colouring-weight} apply the first assertion for the bipartitions with $m\le \frac{2(M-1)}{\Theta} |R|$ to obtain that these contribute at most $\exp\{o(|R |)\}$. For the bipartitions with $m>\frac{2(M-1)}{ \Theta} |R|$ contribute at most
\[
 2^{|R|} (1-p_-)^{2 (M-1) |R|}=\exp\left(- \Omega(|R|) \right),
\]
using again that $\ln n \geq \frac{\ln 2}{M-1} d$.
\end{proof}

\bibliographystyle{plain}
\bibliography{perc}
\end{document}
